\documentclass[notes=show]{beamer}
%%%%%%%%%%%%%%%%%%%%%%%%%%%%%%%%%%%%%%%%%%%%%%%%%%%%%%%%%%%%%%%%%%%%%%%%%%%%%%%%%%%%%%%%%%%%%%%%%%%%%%%%%%%%%%%%%%%%%%%%%%%%%%%%%%%%%%%%%%%%%%%%%%%%%%%%%%%%%%%%%%%%%%%%%%%%%%%%%%%%%%%%%%%%%%%%%%%%%%%%%%%%%%%%%%%%%%%%%%%%%%%%%%%%%%%%%%%%%%%%%%%%%%%%%%%%
\AtBeginSubsection[
  {\frame<beamer>{\frametitle{Outline}   
    \tableofcontents[currentsection,currentsubsection]}}%
]%
{
  \frame<beamer>{ 
    \frametitle{Outline}   
    \tableofcontents[currentsection,currentsubsection]}
}
\usepackage{ulem}
\usepackage{pdflscape}
\setbeamertemplate{caption}{\insertcaption}
\usepackage{color,xcolor,ucs}
\usepackage{beamerthemesplit}
\usepackage{graphicx}
\usepackage{amsmath}
\usepackage{adjustbox}
\usepackage{pdfpages} 
\definecolor{mintgreen}{rgb}{0.6, 1.0, 0.6}
\makeatletter
\@addtoreset{subfigure}{framenumber}% subfigure counter resets every frame
\makeatother
\newcommand*{\Scale}[2][4]{\scalebox{#1}{$#2$}}%
\newcommand*{\Resize}[2]{\resizebox{#1}{!}{$#2$}}%
\usepackage{amsfonts}
\usepackage{booktabs}
\usepackage{color,soul}
\usepackage{pdfpages}
\setboolean{@twoside}{false}
\usepackage{amsmath}
\usepackage{pdflscape} 
\usepackage{eurosym}
\usepackage{amstext} % for \text
\DeclareRobustCommand{\officialeuro}{%
  \ifmmode\expandafter\text\fi
  {\fontencoding{U}\fontfamily{eurosym}\selectfont e}}
\usepackage[caption = false]{subfig}
\usepackage{appendixnumberbeamer} 
\usepackage{graphicx}
\usepackage{mathpazo}
\usepackage{hyperref}
\usepackage{multimedia}
\usepackage{graphicx}
\usepackage{multirow}
\usepackage{subfig}
\usepackage{verbatim}
\usepackage{booktabs, multicol, multirow}
\setcounter{MaxMatrixCols}{10}
\input{Macros.tex}
\AtBeginSection[]
{
  \begin{frame}[noframenumbering]
    \frametitle{Outline for section \thesection}
    \tableofcontents[currentsection]
  \end{frame}
}


\usepackage{xcolor,soul}
\renewcommand<>{\hl}[1]{\only#2{\beameroriginal{\hl}}{#1}}

% https://tex.stackexchange.com/questions/41683/why-is-it-that-coloring-in-soul-in-beamer-is-not-visible
\makeatletter
\newcommand\SoulColor{%
  \let\set@color\beamerorig@set@color
  \let\reset@color\beamerorig@reset@color}
\makeatother
\SoulColor
\AtBeginSection{\frame{\sectionpage}}

\newsavebox\mybox
\newenvironment{aquote}[1]
  {\savebox\mybox{#1}\begin{quote}\openautoquote\hspace*{-.7ex}}
  {\unskip\closeautoquote\vspace*{1mm}\signed{\usebox\mybox}\end{quote}}

\newtheorem{proposition}[theorem]{Proposition}
\newtheorem{Assumption}[theorem]{Assumption}
\newenvironment{stepenumerate}{\begin{enumerate}[<+->]}{\end{enumerate}}
\newenvironment{stepitemize}{\begin{itemize}[<+->]}{\end{itemize} }
\newenvironment{stepenumeratewithalert}{\begin{enumerate}[<+-| alert@+>]}{\end{enumerate}}
\newenvironment{stepitemizewithalert}{\begin{itemize}[<+-| alert@+>]}{\end{itemize} }
\usetheme{Madrid}
\makeatletter
\setbeamertemplate{footline}
{
  \leavevmode%
  \hbox{%
  \begin{beamercolorbox}[wd=.333333\paperwidth,ht=2.25ex,dp=1ex,center]{author in head/foot}%
    \usebeamerfont{author in head/foot}\insertshortauthor~~\beamer@ifempty{\insertshortinstitute}{}{(\insertshortinstitute)}
  \end{beamercolorbox}%
  \begin{beamercolorbox}[wd=.333333\paperwidth,ht=2.25ex,dp=1ex,center]{title in head/foot}%
    \usebeamerfont{title in head/foot}\insertshorttitle
  \end{beamercolorbox}%
  \begin{beamercolorbox}[wd=.333333\paperwidth,ht=2.25ex,dp=1ex,right]{date in head/foot}%
    \usebeamerfont{date in head/foot}\insertshortdate{}\hspace*{2em}
%    \insertframenumber{} / \inserttotalframenumber\hspace*{2ex} % DELETED
  \end{beamercolorbox}}%
  \vskip0pt%
}
\makeatother
\usepackage{pdfpages}
\begin{document}
	\title[]{\large Wage Inequality Wars		\vspace{0.5cm}
	}
	\author[Raffaele Saggio - Econ 561]{Raffaele Saggio}
	\institute[]{UBC}
	\date{\today \\}
	\maketitle

%%%%%%%%%%%%%%%%%%%%%%%%%	
%%%%%%%%%%%%%%%%%%%%%%%%%
%%%%%%%%%%%%%%%%%%%%%%%%%
%%%%%%%%%%%%%%%%%%%%%%%%%
\setbeamertemplate{footline}[frame number]{}

%gets rid of bottom navigation symbols
\setbeamertemplate{navigation symbols}{}

%gets rid of footer
%will override 'frame number' instruction above
%comment out to revert to previous/default definitions
\setbeamertemplate{footline}{}
\begin{frame}
\frametitle{Today}
\begin{itemize}
\item Today we will review several papers that have used a ``neo-classical" paradigm to explain the rise of wage inequality in the United States that started approximately in 1980.
\bigskip
\item We will start by analyzing how an (exogenous?) increase in the supply of a particular group of workers influence the wages of other workers. 
\bigskip
\item We will then move into the so called SBTC literature. This literature argues that increases in wage inequality during the 1980 is driven by increases in technology that have ``pushed up the price" of high skilled workers. 
\bigskip
\item The next lecture will examine papers that have criticized the SBTC change and analyzed other factors that might have played a role to explain the rise in inequality.  
\end{itemize}
\end{frame}
%%
\begin{frame}
\frametitle{Acemoglu, Autor and Lyle (2004)}
\begin{figure}[htb]
\centering
\includegraphics[width=1.1\textwidth]{aux/aal2.png}
\end{figure}
\end{frame}
%%%
\begin{frame}
\frametitle{Acemoglu, Autor and Lyle (2004)}
\begin{figure}[htb]
\centering
\includegraphics[width=0.95\textwidth]{aux/aal1}
\end{figure}
\end{frame}
%%
\begin{frame}
\frametitle{Acemoglu, Autor and Lyle (2004)}
\begin{figure}[htb]
\centering
\includegraphics[width=1\textwidth]{aux/aal3}
\end{figure}
\end{frame}
%%
\begin{frame}
\frametitle{Acemoglu, Autor and Lyle (2004)}
\begin{itemize}
\item WWII pushed a lot of women into the labor force
\bigskip
\item What this (plausibly exogenous) supply shock implied for the earnings of men (and other women)?
\bigskip
\item Exploit state level variation in WWII effects on female LFP.
\begin{itemize}
\item Massachusetts: 55 percent of men between the ages of 18 and 44 left the labor market to serve in the war.
\item Georgia: $\approx$ 40-45 percent.
\end{itemize}
\pause 
\medskip
\item Why these differences?
\begin{itemize}
\item Exemptions for farmers; Age composition; Ethic and occupational structures.
\item Country of origin composition (Italy, German, Japanese).
\item Idiosyncratic differences in the behavior of local draft boards.
\end{itemize}
\end{itemize}
\end{frame}
%%
%%
\begin{frame}
\frametitle{Acemoglu, Autor and Lyle (2004)}
\begin{figure}[htb]
\centering
\includegraphics[width=1\textwidth]{aux/aal4}
\end{figure}
\end{frame}
%%%
\begin{frame}
\frametitle{Acemoglu, Autor and Lyle (2004)}
\begin{figure}[htb]
\centering
\includegraphics[width=1\textwidth]{aux/placeboAAL}
\end{figure}
\end{frame}
%%%
\begin{frame}[shrink=10]
\frametitle{Model}
\begin{itemize}
\item Think of how wages for men and female are set in a competitive market (State) characterized by different levels of men and women. %price is endogenous; %supply are exogenous. 
\pause
\item Three factor nested CES
\be
Y_{t}=A_{t}K_{t}^{\alpha}\left[(1-\lambda)(B_{t}^{M}M_{t})^{\rho}+\lambda(B_{t}^{F}F_{t})^{\rho}\right]^{\frac{1-\alpha}{\rho}}
\ee
\begin{itemize}
\item $\rho\leq 1$
\item $A_{t}\equiv$ neutral productivity.
\item $B_{t}^{G}$: biased technical shifter for gender $G\in\{M;F\}$.
\item Elasticity of substitution between L and K?
\pause
\item Elasticity of substitution b/w M and F: \pause$\sigma_{MF}\equiv\dfrac{1}{1-\rho} $
\end{itemize}
\medskip
\bigskip
\item Competitive paradigm implies that
\be
w_{t}^{F} = (1-\alpha)\lambda B_{t}^{F}A_{t}K_{t}^{\alpha}(B_{t}^{F}F_{t})^{-\alpha}\left[(1-\lambda)\left(\dfrac{B_{t}^{M}M_{t}}{B_{t}^{F}F_{t}}\right)^{\rho}+\lambda\right]^{\frac{(1-\alpha-\rho)}{\rho}}
\ee
\be
w_{t}^{M} = (1-\alpha)(1-\lambda) B_{t}^{M}A_{t}K_{t}^{\alpha}(B_{t}^{M}M_{t})^{-\alpha}\left[(1-\lambda)+\lambda\left(\dfrac{B_{t}^{F}F_{t}}{B_{t}^{M}M_{t}}\right)^{\rho}\right]^{\frac{(1-\alpha-\rho)}{\rho}}
\ee
\end{itemize}
\end{frame}
%%
\begin{frame}
\frametitle{Derivations}
\begin{itemize}
\item Want: ``short run general equilibrium elasticities" (\underline{fixed capital and supply of men})
\be
\dfrac{\partial \log w_{t}^{F}}{\partial \log F_{t}} \equiv \dfrac{1}{\sigma_{F}} = -(1-s_{t}^{m})\alpha -s_{t}^{m}\dfrac{1}{\sigma_{MF}}
\ee\
where $s_{t}^{m}=\dfrac{w_{t}^{M}M_{t}}{w_{t}^{M}M_{t}+w_{t}^{F}F_{t}}$.
\be
\dfrac{\partial \log w_{t}^{M}}{\partial \log F_{t}} \equiv \dfrac{1}{\sigma_{M}} = -(1-s_{t}^{m})\alpha + (1-s_{t}^{m})\dfrac{1}{\sigma_{MF}}
\ee
\item Case 1: $\sigma_{MF}=\infty$ $\rightarrow$ male and female wages both have elasticity $-(1-s_{t}^{m})\alpha$. Intuition?
\item Case 2: $\alpha=0$ additional females raises men wages. Intuition?
\medskip
\pause
\item Anything weird about these calculations? 
\end{itemize}
\end{frame}
\begin{frame}[shrink=10]
\frametitle{Adding quality}
\begin{itemize}
\item Expand the production function as follows now
\be
Y_{t}=A_{t}K_{t}^{\alpha}\left\{(B_{t}^{L})^{\zeta}+\underbrace{\left[(B_{t}^{F}F_{t})^{\mu}+(B_{t}^{H}H_{t})^{\mu}\right]^{\frac{\zeta}{\mu}}}_{N_{t}}\right\}^{\frac{1-\alpha}{\zeta}}
\ee
\begin{itemize}
\item $\sigma_{FH}=(1-\mu)^{-1}$; $\sigma_{FN}=(1-\zeta)^{-1}$.
\medskip
\item When $\zeta > \mu$ female labor competes more with low skill male labor than high skill male labor.
\pause
\end{itemize}
\item We have now the technology to understand the implications of $\uparrow$ $F_{t}$ on the male skill premium:
\be
\omega_{t}=\dfrac{w_{t}^{H}}{w_{t}^{L}} = \dfrac{B_{t}^{H}}{B_{t}^{L}}\dfrac{\left (B_{t}^{H}H_{t}\right)^{\mu-1}\left[(B_{t}^{F}F_{t})^{\mu}(B_{t}^{H}H_{t})^{\mu}\right]^{\frac{\zeta-\mu}{\mu}}}{(B_{t}^{L}L_{t})^{\zeta-1}}
\ee
\item This implies that
\be
\sign(\dfrac{\partial \log \omega_{t}}{\partial \log F_{t}}) = \sign(\zeta-\mu)
\ee
\item An increase in labor supply increases wage inequality when women compete more with low skill men than with high skill men $(\zeta>\mu)$
\end{itemize}
\end{frame}
%%
\begin{frame}[shrink=10]
\frametitle{Implementation}
\begin{itemize}
\item Data: Microdata from decennial CENSUS (IPUMS) - 1940-1950; 47 States.
\be
\log w_{ist} = \delta_{s} + \gamma_{1950} + f_{i}+X_{ist}'\beta_{t}^{g}+\xi\log\left(\dfrac{F_{st}}{M_{st}}\right)+\eta f_{i}\log\left(\dfrac{F_{st}}{M_{st}}\right) + u_{ist}
\ee
\item $w_{ist}\equiv$ weekly earnings.
\medskip
\item $X_{ist}'\beta_{t}^{g}$ gender specific observables.
\item $\dfrac{F_{st}}{M_{st}}$: ratio of female to male labor supply (expressed in weeks).
\pause
\bigskip
\bigskip
\item Model disciplines interpretation of $\eta$ and $\xi$
\item $\eta=\frac{1}{\sigma_{MF}}$;  $\eta+\xi=1/\sigma_{F}$.
\medskip
\item Theoretical restriction
\be
\dfrac{1}{\sigma_{F}} = -(1-s_{t}^{M})\alpha-(s_{t}^{M}\sigma_{MF}^{-1})
\ee
\item Testing under "calibration": $\alpha=0.33$;  $s_{t}^{M}=0.805$
\pause
\item What happened to capital?
\item What is a ``market" here?
\item Would you actually run this regression?
\end{itemize}
\end{frame}
%%
\begin{frame}
\frametitle{Can reject perfect substitutes and model restriction holds}
\begin{figure}[htb]
\centering
\includegraphics[width=1\textwidth]{aux/aal5}
\end{figure}
\end{frame}
%%
\begin{frame}[shrink=20]
\frametitle{Another estimating equation}
\be
\nn
\log w_{ist}^{M} = \delta_{s} + \gamma_{1950} + c_{i}+X_{ist}'\beta_{t}^{e}+\pi\log\left(\dfrac{F_{st}}{M_{st}}\right)+\theta c_{i}\log\left(\dfrac{F_{st}}{M_{st}}\right) + \nu\log\left(\dfrac{C_{st}^{m}}{H_{st}^{m}}\right)+ r_{ist}
\ee
\begin{itemize}
\item Now just focus on earnings of men: college and high school graduates. 
\bigskip
\item $\sigma_{FH}=\frac{1}{\pi} \rightarrow$ cross elasticity b/w females and male high school graduates.
\bigskip
\item $\sigma_{FC}=\frac{1}{\pi+\theta}\rightarrow$ cross elasticity b/w females and male college graduates.
\bigskip
\item If $\frac{\pi+\theta}{\pi}=\dfrac{\sigma_{FH}}{\sigma_{FC}}<1\rightarrow$ female labor supply has greater impact on high school graduates.
\end{itemize}
\end{frame}
%%
\begin{frame}
\frametitle{Women increased inequality among men}
\begin{figure}[htb]
\centering
\includegraphics[width=0.9\textwidth]{aux/aal6}
\end{figure}
\end{frame}
%%
\begin{frame}
\frametitle{Ratio of cross elasticity b/w 0.4-0.6 but imprecisely estimated}
\begin{figure}[htb]
\centering
\includegraphics[width=0.9\textwidth]{aux/aal6}
\end{figure}
\end{frame}
%%%%%%%%%%%%%%%%%%%%%%%%%%%%%%
\begin{frame}
\frametitle{Katz and Murphy 1992}
\begin{itemize}
\item They document massive growth in various measures of wage inequality in the late 70's  / early 80's using March CPS file.
\bigskip
\pause
\item Let's open a parenthesis here: Every labor economist must know what March CPS file actually means. 
\bigskip
\item CPS stands for the current population survey. It contains two key sources
\bigskip
\begin{itemize}
\item March CPS: retrospective questions on annual earnings, weeks worked and usual hours per week for the year prior to the interview $\rightarrow$ typically used to study annual earnings, weeks worked and ``usual" hours per week in the year prior to the interview
\bigskip 
\item May Outgoing Rotation Group: point-in-time estimates of weekly and hourly wages in the week prior to the interview $\rightarrow$ typically used to study hourly wages
\bigskip
\item Some inconsistencies over the years. Lots of very smart people spent a significant amount of their time harmonizing these datasets (Autor, Katz and Kearney, 2008). That is just part of the job.
\end{itemize}
\end{itemize}
\end{frame}
%%
\begin{frame}
\frametitle{Starting point}
\begin{figure}[htb]
\centering
\includegraphics[width=1\textwidth]{aux/km1}
\end{figure}
Residualized wage means: ``residuals from a regression of log weekly earnings on a quartic in experience fully interacted with sex and four education level dummies and linear terms in education within these categories"
\end{frame}
%%
\begin{frame}
\frametitle{Starting point}
\begin{figure}[htb]
\centering
\includegraphics[width=1\textwidth]{aux/km00}
\end{figure}
\end{frame}
%%
\begin{frame}
\frametitle{Trends in Supply}
\begin{figure}[htb]
\centering
\includegraphics[width=1\textwidth]{aux/km000}
\end{figure}
\end{frame}%%
\begin{frame}
\frametitle{KM Idea}
\begin{figure}[htb]
\centering
\includegraphics[width=1\textwidth]{aux/km_idea}
\end{figure}
\end{frame}
\begin{frame}[shrink=30]
\frametitle{Are these trends supply driven?}
\begin{itemize}
\item Multi-factor framework. Let $X_{t}$ denote all labor inputs. 
\medskip
\item These inputs obey a labor demand law
\be
X_{t} = D(\underbrace{W_{t}}_{\text{Vector of Wages}},\underbrace{Z_{t}}_{\text{Vector of Demand Shifters}})
\ee 
\item Total differentiation implies 
\be
dX_{t} = D_{w}dW_{t}+D_{z}dZ_{t}
\ee
\item Premultiply by $dW_{t}'$ we get
\be
dW_{t}^{'}(dX_{t}-D_{z}dZ_{t}) = dW_{t}'D_{w}dW_{t}
\ee
\bigskip
\item Can we sign the right-hand-side?
\pause
\be
dW_{t}^{'}(\underbrace{dX_{t}-D_{z}dZ_{t}}_{\text{``Net" Change in Inputs}}) = dW_{t}'D_{w}dW_{t}\leq0
\ee
\bigskip
\item Now suppose demand shifters did not change over time ($dZ_{t}=0$), then we have the implication that for any time period $t$ and $\tau$
\be
(W_{t}-W_{\tau})(X_{t}-X_{\tau})\leq 0
\ee
\pause
\item KM evaluate to what extent changes in \underline{relative} wage change is explained by changes in \underline{relative} supplies.
\end{itemize}
\end{frame}
%%
%%
\begin{frame}
\frametitle{Lots of wrong signed estimates, especially in 1985}
\begin{figure}[htb]
\centering
\includegraphics[width=1\textwidth]{aux/km2}
\end{figure}
\end{frame}
%%
\begin{frame}[shrink=10]
\frametitle{Closing the Circle}
\begin{itemize}
\item Conclusion: something must have happened to the \underline{relative} demand for skill $\rightarrow$ Skill biased technological change.
\pause
\medskip
\item Use a CES production function
\be
F_{t}(H_{t},L_{t})= (A_{H_{t}}H_{t}^{\rho}+A_{L_{t}}L_{t}^{\rho})^{\frac{1}{\rho}}
\ee
\item Remember from previous lecture that $\sigma=\dfrac{1}{1-\rho}$
\item $(A_{H_{t}},A_{L_{t}})\equiv$ technology parameters that measure the productivity of high and low skilled labor.
\item In perfectly competitive markets $P=MPL$, therefore
\be
\omega_{t} \equiv \dfrac{w_{H_{t}}}{w_{L_{t}}} = \dfrac{A_{H_{t}}}{A_{L_{t}}}\left(\dfrac{H_{t}}{L_{t}}\right)^{\rho-1}
\ee
\item This implies that the wage gap between high and low skilled workers is given by 
\be
\log \omega_{t} = D_{t}-\dfrac{1}{\sigma}\log\dfrac{H_{t}}{L_{t}}
\ee
where $D_{t}=\log(\frac{A_{Ht}}{A_{Lt}})$
\medskip
\item KM take this model to the data parametrizing $D_{t}$ as a linear trends
\pause
\medskip
\item Wait, aren't we regressing price on quantity?(answer: yes).
\end{itemize}
\end{frame}
%
\begin{frame}
\frametitle{$\sigma=1.4$ is the magical number!}
\centering
\includegraphics[width=1\textwidth]{aux/km1E}
\end{frame}
%
\begin{frame}
\frametitle{Wrap up}
\begin{itemize}
\item KM is the triumph of the supply and demand paradigm.
\medskip
\item Incredibly influential (7.5K GS cites)
\medskip
\item Skill biased technological change is important to explain changes in wage inequality.
\pause
\bigskip
\bigskip
\item Clearly, there are several limitations
\bigskip
\begin{itemize}
\item Demand is unobserved (``name the residual").
\medskip
\item No mentioning of labor market institutions.
\medskip
\item What about imperfect competition?  
\end{itemize}
\end{itemize}
\end{frame}
%%
\begin{frame}
\frametitle{Bound and Johnson (1992)}
\begin{itemize}
\item Demand is clearly important. 
\bigskip
\item Use industry / demography structure to quantitatively decompose wage gaps into components attributable to Supply / Demand / Institutions.
\bigskip
\item Key distinction: other existing theories are not ruled out informally (like in KM) but all alternative theories are encompassed within the same approach! 
\end{itemize}
\end{frame}
%%
\begin{frame}[shrink=30]
\frametitle{Conceptual Framework}
\begin{itemize}
\item Production in industry $j$ is a function of $I$ demographic groups (age $\times$ gender $\times$ education)
\be
Q_{j}= a_{j}\left[\sum_{i=1}\delta_{ij}\left(b_{ij}N_{ij}\right)^{\frac{\sigma-1}{\sigma}}\right]^{\frac{\sigma}{\sigma-1}}
\ee
where $a_{j}$ is a ``neutral" shifter in a given industry; $b_{ij}$ is a skilled biased shifter; $\delta_{ij}: ``institutions"$.
\bigskip
\item Industry product demand obeys
\be
\dfrac{Q_{j}}{Q_{r}} = \theta_{j}P_{j}^{-\epsilon}
\ee
\bigskip
\item Hire until marginal revenue product equals the wage $\rightarrow$ (competitive) wage changes are given by
\be
\log d W_{ic} = \left(1-\frac{1}{\sigma}\right)d \log \underbrace{b_{i}}_{\text{SBTC}} - \dfrac{1}{\sigma}d\log \underbrace{N_{i}}_{\text{Supply}} + \dfrac{1}{\sigma}d\log \underbrace{D_{i}}_{\text{Demand}}
\ee
\begin{itemize}
\item $N_{i} = \sum_{j}N_{ij} \rightarrow$ Total employment
\item $b_{i} = \dfrac{1}{J}\sum_{j}b_{ij} \rightarrow$ technical change for group $i$.
\item $D_{i} = \sum_{j=1}^{J}\delta_{ij}^{\sigma}\left(\dfrac{b_{ij}}{b_{i}}\right)^{\sigma-1}x_{j}$
\item $x_{j}=a_{j}\theta_{j}^{\frac{\sigma}{\epsilon}}Q_{j}^{1-\frac{\sigma}{\epsilon}}$
\end{itemize}
\end{itemize}
\pause
Observed wages may deviate from actual wages because of institutional factors (i.e. unionizations which may vary across industries). 

\end{frame}
%%
\begin{frame}
\frametitle{Identification}
\begin{itemize}
\item Assume that SBTC and Demand follow a linear trend
\bigskip
\item Use second differences to identify $\sigma$:
\be
d^{2} \log W_{i} = constant -\dfrac{1}{\sigma}d^{2}\log N_{i}
\ee
\item BJ estimate $\hat{\sigma}=1.7$. Then use model structure to back up importance of $SBTC$ and Demand factors:
\be
d \log W_{i} = \left(1-\frac{1}{\hat{\sigma}}\right)d \log \underbrace{b_{i}}_{\text{SBTC}} - \dfrac{1}{\hat{\sigma}}d\log \underbrace{N_{i}}_{\text{Supply}} + \dfrac{1}{\hat{\sigma}}d\log \underbrace{D_{i}}_{\text{Demand}}
\ee
\item Industry fixed effects used to infer importance of ``rents" (Krueger and Summer, 1988). 

\end{itemize}
\end{frame}
\begin{frame}
\frametitle{Conclusion: Supply and SBTC extremely important}
\centering
\includegraphics[width=0.9\textwidth]{aux/BJ}
\end{frame}
%%
\begin{frame}
\frametitle{Wrap-up}
\begin{itemize}
\item Conceptual framework to quantify competing forces in explaining wage inequality trends.
\medskip
\item Industry level analysis suggests demand is not that important $\rightarrow$ key is SBTC.
\medskip
\item Lot of important follow up to this work
\medskip
\begin{itemize}
\item Capital skill complementarities (Berman, Bound, Grichiles, 1994; Berman, Bound, Machin,
1998; Krusell, Ohanion, Rios-Rull, Violante, 2000)
\item Correlation with computer usage (Krueger 1993; Autor, Katz and Krueger 1993).
\end{itemize}
\bigskip
\bigskip
\item End of 1990's: near consensus on importance of SBTC.
\pause
\medskip
\item  By mid 2000's: dissensus has emerged. East-West coast war begins 
\bigskip
\begin{itemize}
\item Di Nardo, Fortin and Lemieux, 1996; Card-Lemieux (2001). Card, Di Nardo (2002).
\medskip
\item Institutions more important than originally accounted for.
\end{itemize}
\end{itemize}
\end{frame}
%%%%%%%%%%%%%%%%%%%%%%%%%%%
\end{document}

